%%%%%%%%%%%%%%%%%%%%%%%%%%%%%%%%%%%%%%%%%%%%%%%%%%%%%%%%%%%%%%%%%%%%%%%%%%%
\chapter{The RAG Pipeline}
\label{ch:rag-pipeline}
%%%%%%%%%%%%%%%%%%%%%%%%%%%%%%%%%%%%%%%%%%%%%%%%%%%%%%%%%%%%%%%%%%%%%%%%%%%

Retrieval-Augmented Generation (RAG) is a pattern for answering questions against a corpus
too large to fit in a single LLM prompt. A retrieval system selects the relevant subset; the
LLM uses that subset to generate a grounded answer.

%%%%%%%%%%%%%%%%%%%%%%%%%%%%%%%%%%%%%%%%%%%%%%%%%%%%%%%%%%%%%%%%%%%%%%%%%%%
\section{Two Phases}
%%%%%%%%%%%%%%%%%%%%%%%%%%%%%%%%%%%%%%%%%%%%%%%%%%%%%%%%%%%%%%%%%%%%%%%%%%%

RAG has two distinct phases that run at different times and at different frequencies:

\begin{Verbatim}
INDEXING PHASE (offline, once per document)
  Raw document → extract text → chunk → embed → store in index

QUERY PHASE (online, once per question)
  Question → embed → search index → retrieve chunks → assemble context → generate
\end{Verbatim}

The indexing phase is batch work: slow, expensive, done in advance. The query phase must be
fast: the user is waiting.

%%%%%%%%%%%%%%%%%%%%%%%%%%%%%%%%%%%%%%%%%%%%%%%%%%%%%%%%%%%%%%%%%%%%%%%%%%%
\section{Indexing Phase}
%%%%%%%%%%%%%%%%%%%%%%%%%%%%%%%%%%%%%%%%%%%%%%%%%%%%%%%%%%%%%%%%%%%%%%%%%%%

\subsection{Step 1: Document Ingestion}

Accept the raw document in its source format. Common formats: plain text (\texttt{.txt},
\texttt{.md}), PDF, HTML, DOCX, XLSX, code.

At this step, also extract metadata: source URL or path, title, author, creation date,
document type. This metadata travels with every chunk the document produces.

\textbf{Failure modes:}

\begin{itemize}
    \item PDF with scanned images: text extraction returns nothing. Requires OCR.
    \item PDF with ligature fonts: some extractors corrupt \textit{fi}, \textit{fl} ligatures.
    \item HTML with JavaScript-rendered content: static extraction misses dynamically loaded content.
\end{itemize}

\subsection{Step 2: Text Cleaning}

Before chunking, normalize the extracted text:

\begin{itemize}
    \item Remove boilerplate (headers, footers, page numbers, navigation)
    \item Normalize whitespace and Unicode (ligatures, smart quotes, zero-width spaces)
    \item Optionally repair broken words from PDF column extraction
\end{itemize}

Cleaning quality directly affects embedding quality. Embedding a chunk full of OCR artifacts
produces a vector that does not cluster with clean content.

\subsection{Step 3: Chunking}

Split the cleaned text into chunks sized for the embedding model's context window and the
retrieval precision you need. See Chapter~\ref{ch:chunking} for the full decision tree.

Key decisions at this step:

\begin{itemize}
    \item Chunk size (characters or tokens)
    \item Overlap (0--20\% is typical)
    \item Split strategy (fixed, sentence-boundary, structure-aware, semantic)
    \item Parent-child relationship (store parent for context, index children for retrieval)
\end{itemize}

Output: a list of \texttt{(chunk\_text, chunk\_metadata)} pairs.

\subsection{Step 4: Embedding}

Pass each chunk through the embedding model to produce a vector:

\begin{Verbatim}
chunk_text  → [embedding model] → [0.21, -0.45, 0.03, ...]
\end{Verbatim}

This is the most computationally expensive step. Strategies:

\begin{itemize}
    \item \textbf{Batch embedding:} send multiple chunks per API call (up to 2{,}048 for
          OpenAI). Reduces round-trip overhead by orders of magnitude vs.\ one call per chunk.
    \item \textbf{Local vs.\ API:} local ONNX or GGUF models have zero per-call cost but
          require GPU/CPU capacity. OpenAI API has per-token cost but no infrastructure overhead.
    \item \textbf{Caching:} if re-indexing a partially updated corpus, skip chunks whose text
          has not changed (hash the chunk text to detect unchanged content).
\end{itemize}

\subsection{Step 5: Storing in the Index}

Write each \texttt{(vector, chunk\_text, chunk\_metadata)} into the storage backend.

What gets stored:

\begin{itemize}
    \item The vector (for ANN search)
    \item The chunk text (to return as context to the LLM)
    \item The metadata (for filtering and attribution)
    \item Optionally: a BM25 inverted index over the chunk text (for hybrid search)
\end{itemize}

\subsection{Indexing Pipeline Summary}

\begin{Verbatim}
[raw docs]
    │
    ▼
[extract text]  ←── format-specific parsers
    │
    ▼
[clean text]    ←── normalize, strip boilerplate
    │
    ▼
[chunk]         ←── size, overlap, strategy choice
    │
    ▼
[embed]         ←── model choice, batch size, caching
    │
    ▼
[store]         ←── vector index + optional BM25 index
\end{Verbatim}

%%%%%%%%%%%%%%%%%%%%%%%%%%%%%%%%%%%%%%%%%%%%%%%%%%%%%%%%%%%%%%%%%%%%%%%%%%%
\section{Query Phase}
%%%%%%%%%%%%%%%%%%%%%%%%%%%%%%%%%%%%%%%%%%%%%%%%%%%%%%%%%%%%%%%%%%%%%%%%%%%

\subsection{Step 1: Receive the Query}

The user submits a question or search string. Before retrieval, consider:

\begin{itemize}
    \item \textbf{Query type:} Is this a factual lookup, a synthesis request, or a conversational
          follow-up?
    \item \textbf{Query transformation:} Should the query be rewritten before retrieval?
\end{itemize}

Common query transformation techniques:

\begin{description}
    \item[HyDE (Hypothetical Document Embedding)] Ask the LLM to write a hypothetical document
          that would answer the question, then embed that document instead of the bare query.
          Bridges the vocabulary gap between short queries and long documents.
    \item[Query expansion] Append related terms to improve recall.
    \item[Sub-question decomposition] Split a complex query into simpler sub-questions, retrieve
          for each, then combine.
\end{description}

\subsection{Step 2: Embed the Query}

Embed the query using the same model used to embed the documents. This is non-negotiable---
the query and document vectors must live in the same vector space.

\begin{Verbatim}
query_text → [same embedding model] → query_vector
\end{Verbatim}

If query instructions are used (as with Qwen3-Embedding), wrap the query in the appropriate
instruction template before embedding. See Chapter~\ref{ch:query-instructions}.

\subsection{Step 3: Retrieve Candidates}

\textbf{Pure vector search:}

\begin{Verbatim}
query_vector → ANN search → top-k chunks ranked by vector similarity
\end{Verbatim}

\textbf{Hybrid search (BM25 + vector):}

\begin{Verbatim}
query_text   → BM25 search → lex_results (ranked by keyword relevance)
query_vector → ANN search  → vec_results (ranked by semantic similarity)
lex_results + vec_results   → RRF fusion → merged_results
\end{Verbatim}

Hybrid search consistently outperforms either modality alone, especially on queries with
exact-match terms or where the vocabulary gap between query and document is small.

\subsection{Step 4: Filter}

Apply metadata filters to the candidate set:

\begin{itemize}
    \item Date range (``only documents from the last 6 months'')
    \item Source type (``only from the API documentation'')
    \item Access control (``only documents this user can see'')
\end{itemize}

Filtering is most efficient when done \textbf{during} ANN search (as supported by Qdrant and
MongoDB Atlas) rather than after (as in memvid). Post-filtering on a small top-k can silently
produce fewer results than requested if many candidates are filtered out.

\subsection{Step 5: Rerank (Optional but Recommended)}

Pass the top candidates through a cross-encoder reranker to produce a more precise ordering.

\begin{Verbatim}
(query, candidate_1) → [reranker] → 0.92
(query, candidate_2) → [reranker] → 0.87
(query, candidate_3) → [reranker] → 0.41
\end{Verbatim}

The reranker reads the query and each candidate together, producing a relevance score that
accounts for the full text interaction. Reranking is skipped when: latency is critical, the
corpus is small enough that vector search precision is already high, or no reranker model is
available.

\subsection{Step 6: Assemble Context}

Select the top $N$ chunks and concatenate them into a context block for the LLM prompt.

\textbf{Context budget:} The LLM has a finite context window. A practical rule: reserve 30\%
for the answer, 10\% for system prompt and question, 60\% for retrieved context.

\textbf{Ordering matters:} LLMs attend less reliably to content in the middle of long contexts
(the ``lost in the middle'' problem). Put the most relevant chunks first and last.

\textbf{Attribution:} Include source metadata so the LLM can cite sources:

\begin{Verbatim}
[Source: retrieval-concepts.md, Section: HNSW]
HNSW builds a multi-layer graph where each node is a document vector...
\end{Verbatim}

\textbf{Parent-child expansion:} If using parent-child chunking, swap retrieved child chunks
for their parent chunk at this step to give the LLM richer context.

\subsection{Step 7: Generate the Answer}

Typical prompt structure:

\begin{Verbatim}
SYSTEM:
You are a helpful assistant. Answer using only the provided context.
If the context does not contain the answer, say so.

CONTEXT:
[assembled context block]

USER:
[original question]

ASSISTANT:
\end{Verbatim}

The system prompt should instruct the LLM to use the provided context as the primary source,
acknowledge when the context does not contain the answer, and cite sources when relevant.

%%%%%%%%%%%%%%%%%%%%%%%%%%%%%%%%%%%%%%%%%%%%%%%%%%%%%%%%%%%%%%%%%%%%%%%%%%%
\section{Common Failure Modes}
%%%%%%%%%%%%%%%%%%%%%%%%%%%%%%%%%%%%%%%%%%%%%%%%%%%%%%%%%%%%%%%%%%%%%%%%%%%

\subsection{Retrieval Failures}

\textbf{Symptom: The LLM says ``I don't know'' but the answer is in the corpus.}

Probable causes:

\begin{itemize}
    \item Chunk boundary split the answer across two chunks
    \item Query vocabulary doesn't match document vocabulary (semantic gap); try hybrid search or HyDE
    \item \texttt{top\_k} too small; the right chunk is at rank~12 but only top~5 are retrieved
    \item Metadata filter too aggressive; filtered out the relevant chunk
\end{itemize}

\textbf{Symptom: Retrieved chunks are consistently off-topic.}

Probable causes:

\begin{itemize}
    \item Chunks are too large; the relevant sentence is buried in an irrelevant chunk
    \item Embedding model is not well-suited to the domain
    \item No reranker; vector proximity is imprecise for this query type
\end{itemize}

\subsection{Generation Failures}

\textbf{Symptom: LLM ignores the context and answers from its parametric memory.}

Probable causes:

\begin{itemize}
    \item System prompt does not strongly instruct the model to use context
    \item Context block is too long; relevant chunks are in the ``lost in the middle'' zone
    \item Model temperature too high; try 0.0--0.3 for fact retrieval tasks
\end{itemize}

\subsection{Indexing Failures}

\textbf{Symptom: Re-indexing takes much longer than expected.}

Probable causes:

\begin{itemize}
    \item Embedding every chunk individually instead of batching
    \item No caching of embeddings for unchanged content
\end{itemize}

\textbf{Symptom: Index quality degrades over time as documents are updated.}

Probable causes:

\begin{itemize}
    \item Stale vectors from old document versions still in the index
    \item Backend does not support correct deletions (see Chapter~\ref{ch:backends})
    \item No mechanism to detect and re-embed changed content
\end{itemize}

%%%%%%%%%%%%%%%%%%%%%%%%%%%%%%%%%%%%%%%%%%%%%%%%%%%%%%%%%%%%%%%%%%%%%%%%%%%
\section{This Repo's Stack}
%%%%%%%%%%%%%%%%%%%%%%%%%%%%%%%%%%%%%%%%%%%%%%%%%%%%%%%%%%%%%%%%%%%%%%%%%%%

In this repo, the RAG pipeline maps to the following components:

\begin{center}
\begin{tabular}{ll}
\toprule
\textbf{Pipeline step} & \textbf{Component} \\
\midrule
Document ingestion & memvid readers (PDF, DOCX, XLSX) or direct \texttt{put\_bytes()} \\
Chunking & memvid chunker (1{,}200 char, structure-aware) \\
Embedding & \texttt{llama-server --embedding} serving a GGUF embedding model \\
Storage & memvid \texttt{.mv2} file (vector + BM25 + metadata) \\
Query embedding & Same llama-server endpoint \\
Retrieval & memvid \texttt{ask()} --- hybrid BM25 + vector with RRF \\
Reranking & Not yet integrated (planned: Qwen3-Reranker via llama-server) \\
Context assembly & memvid \texttt{AskResponse.context} field \\
Generation & llama-server chat completions or cloud provider via clawrouter \\
\bottomrule
\end{tabular}
\end{center}

The embedding model and generation model are served as separate llama-server instances.
Open WebUI routes embedding requests to the dedicated embeddings endpoint and generation
requests to the generation endpoint. clawrouter handles routing to local vs.\ cloud providers.
